\documentclass[11pt]{article}

\usepackage[preprint]{acl}

\usepackage{times}
\usepackage{latexsym}

\usepackage[T1]{fontenc}
\usepackage[utf8]{inputenc}

\usepackage{microtype}
\usepackage{inconsolata}
\usepackage{graphicx}

\title{DiffusEdit: Interim Proposal}

\author{
  \textbf{Tatva Agarwal\textsuperscript{*}} \\ 2024117005 \And
  \textbf{Vedant Kulkarni\textsuperscript{*}} \\ 2024101140 \And
  \textbf{Byasani Vedavyas\textsuperscript{*}} \\ 2024101025 \And
  \textbf{Venya Velmurugan\textsuperscript{*}} \\ 2024114002
}

\begin{document}

\maketitle

\renewcommand{\thefootnote}{\fnsymbol{footnote}}
\footnotetext[1]{Equal contribution; authors listed alphabetically.}
\renewcommand{\thefootnote}{\arabic{footnote}}
\setcounter{footnote}{0}

\begin{abstract}
Updating documents by hand with the latest available information is a task that requires a lot of human effort. This project attempts to tackle this problem by identifying stale spans of text in a document given the latest information, masking the tokens that make up these stale spans, and re-generating them using a Diffusion Model. This project is built on top of two similar papers: a pre-print titled Detect, Remask, Repair \citep{zou-etal-2026-detect} and a publication titled FRUIT \citep{iv-etal-2022-fruit}.
\end{abstract}

\section{Introduction}

% https://github.com/tchewik/isanlp_rst
% https://github.com/google-research/language/blob/master/language/fruit/README_PIPELINE.md

\bibliography{citations}

% \appendix

% \section{Example Appendix}
% \label{sec:appendix}

% This is an appendix.

\end{document}
